\documentclass[11pt]{article}
\usepackage[margin=1in]{geometry}
\usepackage{amsmath,amssymb,amsthm,mathtools,booktabs,array}
\usepackage[T1]{fontenc}
\usepackage{lmodern,microtype}
\usepackage[colorlinks=true,linkcolor=blue!55!black,urlcolor=blue!55!black,citecolor=blue!55!black]{hyperref}
\usepackage{xcolor}
\newtheorem{theorem}{Theorem}
\newtheorem{lemma}[theorem]{Lemma}
\newtheorem{corollary}[theorem]{Corollary}
\newtheorem{proposition}[theorem]{Proposition}
\theoremstyle{definition}\newtheorem{definition}[theorem]{Definition}
\newcommand{\Rea}{\operatorname{Re}}
\newcommand{\Imb}{\operatorname{Im}}
\newcommand{\ord}{\operatorname{ord}}
\newcommand{\res}{\operatorname{Res}}
\newcommand{\diag}{\operatorname{diag}}
\newcommand{\GRH}{\mathrm{GRH}}
\newcommand{\thq}{\theta_q}
\newcommand{\psq}{\psi_q}
\newcommand{\ee}{\mathrm{e}}
\title{Finite-scale prime-product inequalities\\and sharp one-sided detection of Dirichlet zeros}
\author{Research note prepared with ChatGPT}
\date{26 November 2026}
\begin{document}
\maketitle
\begin{abstract}
Motivated by the one-sided Landau criterion in the Alp\"oge--Claude proof of the Dirichlet generalized Riemann hypothesis (GRH), we construct integer-coefficient dilation filters of the logarithm of the product of primes congruent to $1$ modulo $q$. The filters remove the main prime-counting term and the prime-square bias. A single one-sided, square-root-plus-epsilon bound on their output is equivalent to GRH for the characters modulo $q$. Two fixed dilation bases, $4$ and $9$, recover the exact supremum of the real parts of their nontrivial zeros as a one-sided growth exponent. Under GRH we further obtain lower bounds for both signs of the normalized oscillation, using a Hermitian matrix of residues at prescribed differences of frequencies. The bounds allow coincidences and multiplicities of zeros and require no linear independence hypothesis. The analytic ingredients are classical; the combined filter, exponent and matrix formulations are proposed research contributions, with priority not established by the limited literature search.
\end{abstract}

\section{Source analysis and scope}
The compact edition of \texttt{lsernau/grh-lean} at commit
\texttt{72fead1} (full hash in the bibliography) was inspected through GitHub \cite{sernau}. Its README reports a successful build, kernel replay and Comparator check of the five stated results, with only the usual three Lean axioms. Its authors expressly distinguish machine acceptance from a human reading of the entire proof. Rubinstein reports independent checks of the original development \cite{rubinstein}. We did not rebuild either development; our mathematical proof below is independent of its geometric machinery except when invoking its final GRH result for the forward consequences. The scope is Dirichlet $L$-functions, not general automorphic or nonabelian Artin $L$-functions.

\begin{center}
\small
\begin{tabular}{@{}>{\raggedright\arraybackslash}p{.42\linewidth}>{\raggedright\arraybackslash}p{.51\linewidth}@{}}
\toprule
Compact edition location & Mathematical role \\
\midrule
\texttt{src/Statements.lean} & Precise Mathlib statements, including nonvanishing in $\Rea s>1/2$ (principal pole excluded). \\
\texttt{src/Main.lean} & Five public theorems; \texttt{LFunction\_ne\_zero\_of\_}\allowbreak\texttt{half\_lt\_re} invokes \texttt{New.grh}. \\
\texttt{GUIDE.md}; \texttt{src/New/Main.lean} & Short and sufficient theta lifts, integral long forms, height and index estimates; their comparison yields \texttt{star\_inequality}. \\
\texttt{src/New/Main.lean} & \texttt{eventually\_primeError\_le\_rpow} turns the main inequality and a window argument into a one-sided remainder estimate. \\
\texttt{src/Lib/Analytic/Landau/}\newline\texttt{PolarLandau.lean} & \texttt{IsPolar.landau}: a sign condition on real residues, together with a one-sided exponential bound, rules out poles further right. \\
\bottomrule
\end{tabular}
\end{center}
The key inequality in the compact code has the shape $mR\le (\log b)^2\mathcal B$, where $m$ counts lifts, $R$ is a remainder and $\mathcal B$ is a scale-dependent error budget. The saving comes from $m\asymp\sqrt{x}$. The final analytic step passes to a larger modulus and uses quadratic twists and a one-sided Landau criterion. We transfer the last technique, not the underlying automorphic argument: we design a different remainder whose Mellin residues already have the requisite sign. The source's proof of GRH and our derivation of new consequences should not be confused.

The Wikipedia article \cite{wiki} discusses the link between critical-strip zeros and prime-counting errors, the von Koch criterion, and other equivalent forms. The classical explicit formula and Landau's lemma also appear in standard references \cite{davenport,mv,dlmf}. Our goal is a test with \emph{no explicit main-term subtraction}, involving only products of primes in one congruence class, even though all characters modulo $q$ are detected.

\section{Statements}
Fix an integer $q\ge1$; $\phi(q)$ is Euler's totient, including $\phi(1)=1$. Interpret every integer as congruent to $1$ modulo $1$. Write
\[
 A_q(x)=\prod_{\substack{p\le x\\p\equiv1\ (q)}}p,\qquad
 \thq(x)=\log A_q(x),\qquad
 \psq(x)=\sum_{\substack{n\le x\\n\equiv1\ (q)}}\Lambda(n).
\]
An empty product is $1$. Let $T_bf(x)=f(bx)$ for $b>1$. For a real polynomial $P(z)=\sum_{j=0}^d c_jz^j$, define
\[
 D_{q,b,P}(x)=P(T_b)\thq(x)=\sum_{j=0}^d c_j\thq(b^j x).
\]
Let $\Theta_q$ denote the supremum of the real parts of the zeros in $0<\Rea s<1$ of all Dirichlet $L$-functions modulo $q$, counted or not counted with multiplicity. The principal character includes the zeta zeros, and the functional equation (with conjugation and passage to the inducing primitive character) gives $1/2\le\Theta_q\le1$. Write $\GRH_q$ for $\Theta_q=1/2$.

\begin{definition}[Admissible filter]
For $b>1$, put $r=\sqrt b$. A nonzero real polynomial $P$ is admissible if
\begin{enumerate}
\item $P(b)=P(r)=0$;
\item $P(z)\ne0$ whenever $r<|z|<b$;
\item $P(t)<0$ for every real $r<t<b$.
\end{enumerate}
\end{definition}
The real sign is essential for a one-sided upper bound. The simplest family we use has integer coefficients when $b$ is an integer square:
\[
 P_{b,k}(z)=(z-b)(z-\sqrt b)^{2k},\qquad k\ge1.
\]
The extra factors can instead be chosen to kill selected critical-line frequencies: multiply $P_{b,k}$ by any finite product of factors $z^2-2r\cos(\tau\log b)z+r^2$. Such a factor is positive on $(r,b)$ and has both roots on $|z|=r$. We allow real coefficients in this more general case.

\begin{theorem}[One-sided filter criterion]\label{thm:criterion}
For every admissible $P$, every $b>1$ and every $1/2<\beta<1$, the bound
\[
 D_{q,b,P}(x)\le C x^\beta\quad(x\ge x_0)
\]
for some $C,x_0$ implies that every Dirichlet $L$-function modulo $q$ is zero-free in $\beta<\Rea s<1$. Conversely, if $\Theta_q\le\beta$, then
\[
 |D_{q,b,P}(x)|=O_{q,b,P}(x^\beta\log^2(2x)).
\]
In particular, $\GRH_q$ is equivalent to the one-sided bound $D_{q,b,P}(x)\le C_\epsilon x^{1/2+\epsilon}$ for every $\epsilon>0$ and all sufficiently large $x$. For integer $b$ the equivalence is unchanged if one tests only positive integers $x=n$.
\end{theorem}

\begin{theorem}[Exact one-sided exponent]\label{thm:exponent}
Fix any $k\ge1$ and put $D_b=D_{q,b,P_{b,k}}$, $b=4,9$. Define
\[
 U(n)=\max\{0,D_4(n),D_9(n)\},\qquad
 a^+=\inf\{a\ge0: U(n)=O(n^a)\}.
\]
Unconditionally, $a^+=\Theta_q$. In particular the now machine-verified Dirichlet GRH result implies $a^+=1/2$ for every $q$.
\end{theorem}

\begin{theorem}[Quantitative oscillation under GRH]\label{thm:osc}
Assume $\GRH_q$. Fix any admissible filter. For real $\tau$ put
\[
 r(\tau)=-\frac{P(b^{1/2+i\tau})}{\phi(q)(1/2+i\tau)}\sum_{\chi\bmod q}\ord_{1/2+i\tau}L(s,\chi),
\]
where order means zero multiplicity (zero at a nonzero value of $L$). In particular $r(0)=0$. For any distinct real $\tau_0,\ldots,\tau_m$, the matrix
\[
 R_{ij}=r(\tau_i-\tau_j)\quad(0\le i,j\le m)
\]
is Hermitian, and
\begin{align*}
 \limsup_{x\to\infty}\frac{D_{q,b,P}(x)}{\sqrt x}&\ge\lambda_{\max}(R),\\
 \liminf_{x\to\infty}\frac{D_{q,b,P}(x)}{\sqrt x}&\le\lambda_{\min}(R).
\end{align*}
In particular each surviving zero $\rho=1/2+i\gamma$, $\gamma\ne0$, gives both a positive and negative oscillation of amplitude at least
\[
 \frac{|P(b^\rho)|}{\phi(q)|\rho|}\sum_{\chi\bmod q}\ord_\rho L(s,\chi).
\]
For integer $b$ all the limsup and liminf conclusions hold on the integers as well.
\end{theorem}

\begin{corollary}[Explicit integer prime-product comparisons]\label{cor:products}
Under $\GRH_q$, for every $\epsilon>0$ the inequalities
\begin{align*}
 \log\frac{A_q(64n) A_q(4n)^{20}}{A_q(16n)^8 A_q(n)^{16}}&\le C_{q,\epsilon} n^{1/2+\epsilon},\\
 \log\frac{A_q(729n) A_q(9n)^{63}}{A_q(81n)^{15} A_q(n)^{81}}&\le C_{q,\epsilon} n^{1/2+\epsilon}
\end{align*}
hold eventually. In fact either inequality alone, holding for every $\epsilon>0$, is equivalent to $\GRH_q$. For at least one of these two ratios, the logarithm exceeds $c\sqrt n/\phi(q)$ and is less than $-c\sqrt n/\phi(q)$ infinitely often, for a single absolute constant $c>0$ and a choice of base independent of $q$. Thus even after cancellation of the prime-square bias, neither sign can be bounded by $o(\sqrt n)$ for that base.
\end{corollary}

\section{Mellin transform and the prime-square cancellation}
For $\Rea s>1$, absolute convergence, character orthogonality at the identity, and the Euler product give
\begin{equation}\label{eq:mellin}
 M_q(s):=\int_1^\infty\psq(x)x^{-s-1}\,dx
 =\frac{-1}{s\phi(q)}\sum_{\chi\bmod q}\frac{L'}{L}(s,\chi).
\end{equation}
Indeed, integrating the indicator $x\ge n$ gives $n^{-s}/s$; the sum of all characters is $\phi(q)$ at the identity and $0$ at all other residue classes, also at nonunits. For $q=1$ the sum has just the zeta term. Imprimitive Euler factors can have zeros on $\Rea s=0$, but none in the open critical strip. The continuation of \eqref{eq:mellin} to $\Rea s>0$ has at most simple poles. At $1$ the residue is $1/\phi(q)$; at a zero $\rho$ in the strip it is
\begin{equation}\label{eq:residue0}
 -\frac{m_q(\rho)}{\phi(q)\rho},\qquad
 m_q(\rho)=\sum_{\chi\bmod q}\ord_\rho L(s,\chi)>0.
\end{equation}
No zero cancellation is possible in this sum: multiplicities add. This is why we use the identity residue class. The classical nonvanishing on $\Rea s=1$ ensures that there are no other poles on that line.

For $F(x)=P(T_b)\psq(x)$, change variables $u=b^jx$ in the Mellin integral to get
\begin{equation}\label{eq:filtered}
 \mathcal F(s)=\int_1^\infty F(x)x^{-s-1}\,dx
 =P(b^s)M_q(s)-H(s),
\end{equation}
where
\[
 H(s)=\sum_{j=0}^d c_j b^{js}\int_1^{b^j}\psq(u)u^{-s-1}\,du
\]
is entire. The pole at $1$ is removable because $P(b)=0$. At each zero $\rho$ the residue is
\begin{equation}\label{eq:r}
 -\frac{P(b^\rho)m_q(\rho)}{\phi(q)\rho}.
\end{equation}
For $\Rea\rho>1/2$ in the strip this is nonzero by admissibility. At a real zero $\rho\in(1/2,1)$ it is strictly positive. This sign, opposite to that of the unfiltered prime-power error, is what permits an upper rather than a lower bound to detect every zero, including a possible real exceptional zero.

\begin{lemma}[Prime powers]\label{lem:squares}
For fixed $q$, let $h_q=|\{u\in(\mathbb Z/q\mathbb Z)^*:u^2=1\}|$. Then
\[
 \psq(x)-\thq(x)=\frac{h_q}{\phi(q)}\sqrt{x}+o_q(\sqrt x).
\]
Consequently $F(x)-D_{q,b,P}(x)=o_{q,b,P}(\sqrt x)$.
\end{lemma}
\begin{proof}
Separate prime powers of exponent $2$ from higher powers. The former contribute precisely $\sum_{u^2=1}\theta(\sqrt{x};q,u)$, where $\theta(y;q,u)=\sum_{p\le y,\,p\equiv u\ (q)}\log p$. By the classical prime number theorem in fixed arithmetic progressions each summand is $\sqrt{x}/\phi(q)+o_q(\sqrt{x})$. For exponents $m\ge3$ there are at most $\log x/\log2$ possibilities for $m$, and each contributes at most $x^{1/3}\log x$. This is $O(x^{1/3}\log^2 x)=o(\sqrt x)$. Applying $P(T_b)$ to the leading square term multiplies it by $P(\sqrt b)=0$; the finite sum of remainders stays $o(\sqrt x)$.
\end{proof}
This cancellation is useful even under GRH: passing from $\psi$ to $\theta$ without it can move the center of normalized oscillation by a nonzero constant. The extra root at $\sqrt b$ also removes the possible central zero of a Dirichlet $L$-function from the filtered residue.

\section{The one-sided argument}
For completeness we prove the particular Landau principle used here, rather than treating the library theorem as a black box.
\begin{lemma}[One-sided pole exclusion]\label{lem:landau}
Let $F$ be real and locally integrable on $[1,\infty)$, $F(x)=O(x\log^A(2x))$ for some fixed $A\ge0$. Suppose its Mellin integral continues meromorphically to $\Rea s>\beta$ with at most simple poles, where $\beta<1$, and every real pole in this region has nonnegative residue. If $F(x)\le Cx^\beta$ eventually (or almost everywhere eventually), the continuation has no pole in $\Rea s>\beta$.
\end{lemma}
\begin{proof}
Choose $X$ so that $g(x)=Cx^\beta-F(x)\ge0$ almost everywhere for $x\ge X$; enlarge $C$ to be positive if needed. For $\Rea s>1$ its transform is
\[
 G(s)=\int_X^\infty g(x)x^{-s-1}\,dx
 =\frac{CX^{\beta-s}}{s-\beta}-\mathcal F(s)+\int_1^X F(x)x^{-s-1}\,dx.
\]
Let $\alpha$ be the abscissa of convergence of this nonnegative integral on the real axis, with $-\infty$ allowed. Then $\alpha\le1$, and the integral is holomorphic for $\Rea s>\alpha$ (dominated on compact sub-half-planes, also after any finite number of derivatives). By uniqueness of continuation, if $\mathcal F$ had a pole with real part $>\beta$, necessarily $\alpha>\beta$. Indeed such a pole cannot lie strictly to the right of $\alpha$.

At the real point $\alpha>\beta$, the meromorphic continuation of $G$ has either no pole or a simple pole of nonpositive residue. A genuine pole would have strictly negative residue, contradicting $G(\sigma)\ge0$ as real $\sigma\downarrow\alpha$. Hence $G$ is analytic at $\alpha$. We show this is impossible. Choose a disk of radius $\eta>0$ about $\alpha$ on which the continuation is analytic. Set $s_0=\alpha+\eta/4$ and $s_1=\alpha-\eta/4$. The Taylor series centered at $s_0$ converges at $s_1$, since a disk of radius $3\eta/4$ about $s_0$ lies in the original disk. For each $j\ge0$, differentiation under the convergent integral gives
\[
 (-1)^jG^{(j)}(s_0)=\int_X^\infty g(x)x^{-s_0-1}(\log x)^j\,dx.
\]
Take $X\ge1$. All terms in the Taylor series evaluated at $s_1$ are nonnegative, so monotone convergence identifies its finite sum with $\int_X^\infty g(x)x^{-s_1-1}\,dx$. This contradicts $s_1<\alpha$. Thus the alleged pole cannot exist.
\end{proof}
\begin{proof}[Proof of Theorem~\ref{thm:criterion}]
The prime-power function is $O(x\log(2x))$ trivially, so the same is true of $F$. If $D\le Cx^\beta$, Lemma~\ref{lem:squares} gives $F\le C'x^\beta$ eventually since $\beta>1/2$. Equations \eqref{eq:filtered}--\eqref{eq:r} verify every analytic and sign hypothesis of Lemma~\ref{lem:landau}; there are no poles at or to the right of $1$. Every zero with $\beta<\Rea\rho<1$ would survive, contradicting the lemma.

For the converse we use the standard fixed-modulus truncated explicit formula \cite{davenport,mv,rs}. To spell out the estimate, for each primitive character inducing a character modulo $q$, use a truncation of height comparable to $x$ (or take the standard symmetric formula at height $x$ with its endpoint error). Its principal contribution is $\delta_\chi x$, its zero sum is $-\sum_{|\Imb\rho|\le T}x^\rho/\rho$, and the remaining error is $O_q(\log^2(2x))$ at this choice of height. The possible endpoint half-weight changes the sum by $O(\log x)$; imprimitive Euler-factor corrections to the prime sum are also $O_q(\log x)$. The zero-count estimate $N_\chi(T)=O_q(T\log(2T))$ and partial summation give $\sum_{|\Imb\rho|\le T}|\rho|^{-1}=O_q(\log^2(2T))$ (the finitely many small ordinates give a fixed constant). Thus, if every real part is at most $\beta$, orthogonality yields
\[
 \psq(x)=\frac{x}{\phi(q)}+O_q(x^\beta\log^2(2x)).
\]
The $q$-dependent constants may depend on $\beta$ as well, although for $\beta\in[1/2,1]$ the displayed derivation permits the same constant. Applying $P(T_b)$ kills the main term; Lemma~\ref{lem:squares} deals with the replacement by ordinary primes. The equivalence follows by taking arbitrarily small positive $\epsilon$ and using the functional equation for the conjugate inducing primitive characters. Conversely $x^{1/2}\log^2(2x)=O_\epsilon(x^{1/2+\epsilon})$.

Finally if $b$ is an integer and $n\le x<n+1$, then each $\thq(b^jx)-\thq(b^jn)$ is bounded in absolute value by $b^j\log(b^j(n+1))$: at most $b^j$ integers enter that interval. Hence
\begin{equation}\label{eq:integer}
 |D(x)-D(n)|=O_{b,P}(\log(2n)).
\end{equation}
This extends bounds on integers to real $x$ and also transfers normalized limsup and liminf (using $\sqrt{x}/\sqrt n\to1$ and the known upper bound under GRH, or directly an $o(1)$ difference using that upper bound).
\end{proof}

\section{A Hermitian residue bound for oscillations}
\begin{proof}[Proof of Theorem~\ref{thm:osc}]
Under $\GRH_q$, write $V(t)=\ee^{-t/2}D(\ee^t)$, $t\ge0$. The explicit formula estimate gives $V(t)=O_{q,b,P}((1+t)^2)$. Its exponentially damped Fourier averages therefore exist for every $\delta>0$. For any real $\tau$,
\begin{equation}\label{eq:abel}
 \lim_{\delta\downarrow0}\delta\int_0^\infty V(t)\ee^{-(\delta+i\tau)t}\,dt=r(\tau).
\end{equation}
For the prime-power version this is \eqref{eq:filtered}, evaluated at $s=1/2+\delta+i\tau$ and multiplied by $\delta$. The Mellin integral equals its continuation throughout $\Rea s>1/2$ by the GRH estimate, and the meromorphic function has at most a simple pole at the point in question. Lemma~\ref{lem:squares} makes the difference between the two versions $o(1)$ after normalization; its contribution to \eqref{eq:abel} tends to zero by splitting the integral at any fixed sufficiently large $t$. Conjugation of the family of all characters gives $r(-\tau)=\overline{r(\tau)}$; $r(0)=0$ because $P(\sqrt b)=0$.

Let $L^+=\limsup V$. If $L^+=+\infty$ the first desired bound is immediate. Otherwise take any real $C>L^+$ (including arbitrary $C$ if $L^+=-\infty$). For sufficiently large $t$, $C-V(t)\ge0$. For a complex vector $v=(v_0,\ldots,v_m)$ the quadratic form
\[
 \delta\int_0^\infty (C-V(t))\ee^{-\delta t}
 \left|\sum_{j=0}^m v_j\ee^{i\tau_jt}\right|^2\,dt
\]
is nonnegative in the limit: the integral over a fixed initial interval tends to zero. The limit of the form, by \eqref{eq:abel} and distinctness of the frequencies, is $v^*(CI-R)v$. Thus $CI-R$ is positive semidefinite and $C\ge\lambda_{\max}(R)$. Let $C$ decrease to $L^+$; the case $L^+=-\infty$ is thereby excluded. Similarly, if $C<\liminf V$, use $V-C\ge0$ eventually to obtain $R-CI\succeq0$, hence $\liminf V\le\lambda_{\min}(R)$ (the infinite cases follow in the same manner). The $2\times2$ choice $\{0,\gamma\}$ has eigenvalues $\pm|r(\gamma)|$, proving the individual zero estimate. Equation \eqref{eq:integer} proves the integer assertion.
\end{proof}

\begin{corollary}[Cluster enhancement]\label{cor:cluster}
If $\gamma_1,\ldots,\gamma_m$ are distinct positive ordinates with $r(\gamma_j-\gamma_i)=0$ for $i\ne j$, then the upper and lower normalized limits are at least $B$ and at most $-B$, respectively, where
\[
 B=\left(\sum_{j=1}^m |r(\gamma_j)|^2\right)^{1/2}.
\]
It suffices that all pairwise differences are below the first positive ordinate of any zero in the finite family of $L$-functions in question. The matrix on $0,\gamma_1,\ldots,\gamma_m$ has eigenvalues $B,-B$ and $0$ (with the evident degeneracy when $B=0$).
\end{corollary}
\begin{proof}
The nonzero entries lie in the first row and column, and the resulting rank-at-most-two matrix has the stated spectrum, by restricting to the span of the first coordinate and the vector of residues. Apply Theorem~\ref{thm:osc}.
\end{proof}
No independence of ordinates, simplicity of zeros, or existence of a limiting distribution is required. This complements the classical study of prime races via limiting distributions \cite{rs}. The bounds are lower bounds, not claims that the limsup or liminf is finite or that the matrix estimate is attained. Arbitrarily many prescribed frequencies may be used; the matrix is finite in every application.

\section{Sharpness and the integer examples}
\begin{proof}[Proof of Theorem~\ref{thm:exponent}]
If $\Theta_q<1$, the truncated-formula argument with exponent $\Theta_q$ gives $U(n)=O_{q,k}(n^{\Theta_q}\log^2(2n))$. If $\Theta_q=1$, the trivial estimate $U(n)=O_{q,k}(n\log(2n))$ suffices. Hence $a^+\le\Theta_q$. If $\Theta_q>1/2$ and $a^+<\Theta_q$, then there is an exponent $a<\Theta_q$ with $U(n)=O(n^a)$. Choose $\max(a,1/2)<\beta<\Theta_q$. There exists a zero with real part greater than $\beta$ by the definition of supremum. Theorem~\ref{thm:criterion} with either $b=4$ or $9$ contradicts this, so $a^+\ge\Theta_q$.

It remains to handle $\Theta_q=1/2$. Choose once and for all any nonreal zeta zero $\rho=1/2+i\gamma$ with $\gamma>0$ (existence is classical, and its real part follows here from $\Theta_q=1/2$). It is a zero of the principal-character $L$-function for every $q$, because the finite Euler factors are nonzero at $\Rea s=1/2$. For $P_{b,k}$ on $|z|=\sqrt b$, the only root is $\sqrt b$ itself. Thus $P_{b,k}(b^\rho)=0$ iff $\gamma\log b\in2\pi\mathbb Z$. Both $b=4,9$ cannot have this property: the ratio $\log4/\log9=\log2/\log3$ is irrational, by unique prime factorization. At least one base, independent of $q$ and $k$, survives this zero. Set
\[
 B_k=\max_{b=4,9}\frac{|P_{b,k}(b^\rho)|}{|\rho|}>0.
\]
Theorem~\ref{thm:osc} gives a limsup at least $B_k/\phi(q)$ and a liminf at most $-B_k/\phi(q)$ for a maximizing base. The former gives infinitely many integers with $U(n)>B_k\sqrt n/(2\phi(q))$, so $a^+\ge1/2$. This also proves the two-sign assertion.
\end{proof}
\begin{proof}[Proof of Corollary~\ref{cor:products}]
Expand $P_{4,1}=z^3-8z^2+20z-16$ and $P_{9,1}=z^3-15z^2+63z-81$. Since $\thq=\log A_q$, the two displayed ratios are exactly the filter outputs after taking logarithms. Apply Theorem~\ref{thm:criterion}, and take $c=B_1/2$ in the last proof. No numerical information on the first zeta zero is needed for strict positivity of $c$.
\end{proof}
The use of two bases is a simple exact anti-aliasing argument: one does not need any numerical certification of a critical ordinate modulo $2\pi/\log b$. For any single base one already has the GRH criterion, even if it happens to cancel some critical-line poles. For the two-base result, the conclusion $a^+=\Theta_q$ is unconditional; identifying it with $1/2$ uses the externally established GRH theorem.

\section{Reproducibility, limitations, and relation to prior work}
The accompanying Python program expands the integer filters exactly, verifies the two vanishing moments and higher root multiplicities using integer arithmetic, and computes the prime-product log ratios using a sieve and prefix sums. It also computes illustrative residue amplitudes using the first two zeta zeros from \texttt{mpmath} (not certified enclosures). For example, with $q=1$ and $k=1$, the approximate cluster bounds from the first two zeta ordinates are $3.476$ for $b=4$ and $15.095$ for $b=9$ if those two zeros have multiplicity one and their difference is not an ordinate. These conditions and numerical values are not independently certified here. The numerical example is a sanity check, not part of the proof. The CSV and text log record the actual run. The proof of the lower bounds needs no numerical zero location. All statements are for fixed $q$ (no assertion of uniform constants for growing moduli, except the explicit lower-bound factor $1/\phi(q)$). The choice of $k$, $b$ and the filter is also fixed as $x\to\infty$.

The derivation uses standard explicit formulas, the PNT in arithmetic progressions, elementary character orthogonality, and positivity of Laplace transforms. Classical one-sided RH criteria and prime-race analysis precede this work. Searches for combinations of one-sided Chebyshev criteria, prime products, dilation and finite differences returned standard RH equivalences and explicit formulas, but did not identify the precise combination of square-bias annihilation, a two-base sharp exponent and the residue-difference matrix in this note. This is evidence only for a proposed novel formulation, not a proof of historical priority. No additional major open conjecture is claimed resolved. The external GRH proof is not reproduced here, and the new results are human-readable mathematical proofs, not independently Lean-kernel-checked theorems. A separate audit note lists the source locations and a possible formalization path.

\begin{thebibliography}{9}
\bibitem{sernau} L. Sernau et al., \emph{GRH for Dirichlet L-functions, in Lean 4}, compact edition, 25 Nov. 2026, commit \texttt{72fead1}. \url{https://github.com/lsernau/grh-lean/tree/72fead102f663d5c7a407964a13dc0cf9d977207/grh-compact}.
\bibitem{wiki} Wikipedia, \emph{Riemann hypothesis}, sections on consequences, equivalent criteria and zero locations, accessed 26 Nov. 2026. \url{https://en.wikipedia.org/wiki/Riemann_hypothesis}.
\bibitem{rubinstein} M. Rubinstein, \emph{Alp\"oge/Claude proof of the Generalized Riemann Hypothesis}, verification and working exposition, accessed 26 Nov. 2026. \url{https://www.math.uwaterloo.ca/~mrubinst/alpoge_claude_grh/}.
\bibitem{davenport} H. Davenport, \emph{Multiplicative Number Theory}, 3rd ed., revised by H. L. Montgomery, Springer, 2000 (explicit formulas for Dirichlet characters, prime number theorem in progressions).
\bibitem{mv} H. L. Montgomery and R. C. Vaughan, \emph{Multiplicative Number Theory I: Classical Theory}, Cambridge University Press, 2007 (explicit formula, zero counts, and Landau's lemma, Lemma 15.1).
\bibitem{rs} M. Rubinstein and P. Sarnak, Chebyshev's bias, \emph{Experimental Mathematics} \textbf{3} (1994), 173--197, especially \S2.1. \url{https://www.math.uwaterloo.ca/~mrubinst/publications/Chebyshev.pdf}.
\bibitem{dlmf} NIST Digital Library of Mathematical Functions, \S25.16, mathematical applications of zeta functions. \url{https://dlmf.nist.gov/25.16}.
\end{thebibliography}
\end{document}
